\clearpage
Bemerkung: Bei Rotationsschwingungsspektrum fast gleich große $B_{\nu'}$ und $B_{\nu''}$\\
unterschied von $F$-Abhängigkeit\\

\begin{figure}[H]
    \centering
    \def\svgwidth{.3\columnwidth}
    \subimport{./figures}{vl13_abbildung_0.pdf_tex}
    \label{fig:vl13_abbildung_0}
\end{figure}

Jetzt: Aufgrund unterschiedlicher Potentialkurven und $\overline{r}_{\mathrm{e}}' \neq \overline{r}_{\mathrm{e}}''$ sind sehr verschiedene $B_{\nu'}$ und $B_{\nu''}$ möglich.\\

Die Auswahlregeln für $J$ hängen von der Art des elektronischen Übergangs ab. So ist $\Delta J = 0$ möglich falls sich der elektronsiche Drehimpuls ändert, da bei der Emission oder Absorptio eines Photons mit Drehimpuls $\hbar $ der Gesamtdrehimpuls erhalten bleiben muss. Beispiele einiger Übergänge sind
\begin{itemize}
	\item $\Sigma \to \Sigma$-Übergang: $\Delta J = \pm 1$ P-, R-Zweige
	\item $\Sigma \to \Pi$-Übergang: $\Delta J = 0, \pm 1$ P-, R-, Q-Zweige
\end{itemize}
Für einen Übergang gilt
$$
\begin{aligned}
	\Delta E_{\text{Ges}} &= \Delta \left( E_{\text{el}} + E_{\text{vib}} \right) + \Delta E_{\text{rot}} \\
	\overline{\nu} &= \overline{\nu}_{(\nu', \nu'')} + \overline{\nu}_{(J', J'')} \\
	\overline{\nu} &= \overline{\nu}_{(\nu', \nu'')} + B'J'(J'+1) - B''J''(J''+1) \\
\end{aligned}
$$ 
mit dem Bandenursprung $(\nu', \nu'')$. Es gilt für den
\begin{itemize}
	\item P-Zweig: $\Delta J = -1$, $J'' = J' +1$
		$$
		\overline{\nu}_{\text{P}} = \overline{\nu}_{\nu', \nu''} - (B' + B'')(J'+1) + (B' - B'')(J'+ 1)^2 
		$$ 
		mit $J'= 0,1,2,\ldots$.
	\item R-Zweig: $\Delta J = + 1$, $J' = J'' +$
		$$
		\overline{\nu}_{\text{R}} = \overline{\nu}_{\nu', \nu''}+ (B' + B'') (J''+1) + (B' - B'') (J'' +1)^2
		$$
		mit $J'' = 0,1,2,\ldots$.
	\item Q-Zweig: $\Delta J = 0$, $J' = J''$
		$$
		\overline{\nu}_{\text{Q}} = \overline{\nu}_{\nu', \nu''} + (B' - B'') J'' + (B' - B'')J''^2.
		$$ 
		Diese Gleichung beschreibt die Fortrat-Parabeln.
\end{itemize}

\begin{figure}[H]
    \centering
    \incfig{vl13_abbildung_1}
    \label{fig:vl13_abbildung_1}
\end{figure}

Bemerkungen:
\begin{enumerate}
	\item aus dem Charakter des Spektrums ergeben sich Aussagen über die relativen Lagen der $r_{\mathrm{e}}$
	\item aus Spektrum-Parabeln $\to $ $B_{\nu'}, B_{\nu''} \to r_{\mathrm{e}}' , r_{\mathrm{e}}''$
	\item alles im sichtbaren Spektralbereich, auch homonukleare Moleküle
	\item Auswahlregen zum Teil anders
\end{enumerate}

\textbf{e) Gesamtspektrum}\\
$r_{\mathrm{e}}' > r_{\mathrm{e}}'' \implies B' < B''$, Reihenspektrum\\
\textcolor{red}{bandengrafik}\\

\textbf{f. Dissoziation}
Eine Dissoziation ohne Elektronenübergang $\Delta \nu = \pm 1, \pm 2$ \\
 $$
 (AB) + h \nu(\text{IR}) \to  A+B
$$ 
\textbf{geht nicht}, da eine sollche gegen das Franck-Condon-Prinzip verstößen würde.
\begin{figure}[H]
    \centering
    \incfig{vl13_abbildung_3}
    \label{fig:vl13_abbildung_3}
\end{figure}

Eine Dissoziation \textbf{mit} Elektronenübergang, also
$$
(AB) + h \nu \text{(sichtbar)} \to  A+B^{*}
$$ 
ist möglich.
\begin{figure}[H]
    \centering
    \incfig{vl13_abbildung_4}
    \label{fig:vl13_abbildung_4}
\end{figure}

mit der Konvergenzstelle $K$, der Dissoziationsenergie des Grundzustands $D''$, der Dissoziationsenergie des angeregten Zustand $D'$, der Anregungsenergie des bei der Dissoziation entstehenden angeregten Atoms $E_{\text{At}}$ und der molekularen Anregungsenergie $E_{\text{Mol}}$. Diese Konstante wurden experimentell für das Jodmolekül \ce{J2} bestimmt mit $K = \SI{2,48}{\electronvolt}$, $E_{\text{At}} = \SI{0,94}{\electronvolt}$ und $D'' = \SI{1,54}{\electronvolt}$.\\

\textbf{g) Prädissoziation} 

\begin{figure}[H]
    \centering
    \incfig{vl13_abbildung_5}
    \label{fig:vl13_abbildung_5}
\end{figure}

für $K_{\text{b}}< h\nu < K_{\mathrm{a}}$ gibt es 2 Möglichkeiten
\begin{enumerate}
	\item Innerhalb von $a$ kommt es zu Abregung
	\item Beim Übergang von $a \to b$ kommt es zur Dissoziation unter einem umbau der Elektronenhülle
\end{enumerate}
Es kommt zu einer Lebensdauerverkürzung, $\Delta E$ nimmt zu woraus Linienverbreiterung folgt. Hohe Wahrscheinlichkeit für $h\nu = K_{B}$. \\

\textbf{h) Entstehung von kont. Spektren (\ce{H2}-Lampe)}
\begin{figure}[H]
    \centering
    \incfig{vl13_abbildung_6}
    \label{fig:vl13_abbildung_6}
\end{figure}
Anregungen durch Stöße sind beim \ce{H2} erlaubt, wobei Singulett-Triplett übergänge verboten sind. Wasserstofflampen werden oft als UV-Quellen verwendet. Das Prinzip funktioniert analog für UV-Lampen die Edelgasdimere wie \ce{Xe2} oder \ce{Kr2} verwenden.
